\documentclass[10pt,a4paper]{amsart}
\usepackage[margin=19mm]{geometry}
\usepackage{amsmath,amssymb,mathtools}
\usepackage[scr=rsfs,cal=euler]{mathalfa}
\usepackage{xfrac}
\usepackage[unicode,hidelinks,bookmarks=false]{hyperref}
\newcommand{\ie}{i.e.\ }
\newcommand{\cf}{cf.\ }
\newcommand{\resp}{respectively\ }
\newcommand{\Subsection}{\subsection}
\providecommand{\nobreakdash}{\nobreak\hbox{-}}
\setlength{\emergencystretch}{2em}
\multlinegap0pt
\usepackage{amssymb}
\usepackage{float}
\usepackage{mathtools}
\usepackage{bm}
\usepackage{dsfont}
\usepackage[shortlabels]{enumitem}
\usepackage{algorithm}\usepackage{algpseudocode}
\newtheorem{lemma}{Lemma}
\title{When Are Trade-Off Functions Testable from Finite Samples?}
\author{Kaining Shi, Qiaosen Wang, Cong Ma}
\date{Source: arXiv:2605.10774v1 (CC BY 4.0)}
\begin{document}
\maketitle

\noindent\emph{Source and licence.} When Are Trade-Off Functions Testable from Finite Samples?. Version arXiv:2605.10774v1 (CC BY 4.0), distributed by arXiv under the Creative Commons Attribution 4.0 International licence, http://creativecommons.org/licenses/by/4.0/, as recorded in the arXiv licence metadata for this exact version and shown on the abstract page https://arxiv.org/abs/2605.10774v1. Full licence notice: \texttt{licenses/arxiv-2605.10774-CC-BY-4.0.txt}.\par\smallskip

\noindent\textit{Editorial scope.} Counted unit: the article's lemma lem:convex-hull-gcm (source0.tex lines 1961--1975). The article's complete original proof of it is delivered with it (source0.tex lines 1977--2028, one \texttt{proof} environment of 52 lines). No mathematics is written, repaired or re-derived: the statement and the proof are the article's own lines, in the article's own notation, and no retained line is substituted or re-flowed.  No further source-local context is needed: every cross-reference of every retained line resolves inside the two delivered spans.  Nothing was omitted from any retained span, so the package carries no omission record.  No retained line contains a citation, so the wrapper carries no bibliography and no bibliography entry.  No retained line contains a figure environment, an image-inclusion command or drawing code, so no image is delivered and the package declares no asset dependency.  Editorial formatting only, in addition: an AMS \texttt{amsart} preamble that restates the article's own declarations and macros used by the retained lines, namely the environments \texttt{lemma} on separate counters, together with the standard packages those lines need.  The retained environments print in this document as Lemma 1 in order of appearance, whereas the article prints the same items with the numbering of its own full document.  The complete native source of the article as distributed in the arXiv e-print for this exact version is bundled unchanged as \texttt{sources/200318/source0.tex}.
\begin{lemma}[Convex-hull envelope equals the greatest convex minorant]
\label{lem:convex-hull-gcm}
Let $S$ be a class of measurable sets containing $\varnothing$ and $X$, and define
\[
C:=\{(P(A),\,Q(A^{\complement})):\ A\in S\}\subset [0,1]^2.
\]
For $\alpha\in[0,1]$, define
\[
F(\alpha):=\inf_{A\in S:\,P(A)\leqslant \alpha} Q(A^{\complement}),
\qquad
G(\alpha):=\inf\{\beta:\ (\alpha,\beta)\in \operatorname{conv}(C)\}.
\]
Then $G=\operatorname{GCM}(F)$, 
where $\operatorname{GCM}(F)$ denotes the greatest convex minorant of $F$ on $[0,1]$.
\end{lemma}

\begin{proof}
We first show that $G\leqslant F$ pointwise. Fix $\alpha\in[0,1]$ and $A\in S$ with
\[
x:=P(A)\leqslant \alpha,
\qquad
y:=Q(A^{\complement}).
\]
If $x=\alpha$, then $(\alpha,y)\in C\subset \operatorname{conv}(C)$, hence $G(\alpha)\leqslant y$.

If $x<\alpha$, use that $X\in S$, so $(1,0)\in C$. Set
\[
\theta:=\frac{1-\alpha}{1-x}\in[0,1].
\]
Then
\[
\theta(x,y)+(1-\theta)(1,0)=(\alpha,\theta y)\in \operatorname{conv}(C),
\]
so
\[
G(\alpha)\leqslant \theta y\leqslant y.
\]
Taking the infimum over all $A\in S$ with $P(A)\leqslant \alpha$ yields
\[
G(\alpha)\leqslant F(\alpha).
\]

Next we show that $G$ dominates every convex minorant of $F$. Let $g$ be any convex function on
$[0,1]$ such that $g\leqslant F$. For every $A\in S$,
\[
g(P(A))\leqslant F(P(A))\leqslant Q(A^{\complement}),
\]
since $A$ itself is admissible in the definition of $F(P(A))$. Hence
\[
C\subset \operatorname{epi}(g):=\{(u,v)\in[0,1]^2:\ v\geqslant g(u)\}.
\]
Because $g$ is convex, $\operatorname{epi}(g)$ is convex. Therefore
\[
\operatorname{conv}(C)\subset \operatorname{epi}(g).
\]
So for every $(\alpha,\beta)\in \operatorname{conv}(C)$,
\[
\beta\geqslant g(\alpha).
\]
Taking the infimum over such $\beta$ gives
\[
G(\alpha)\geqslant g(\alpha)
\qquad\forall \alpha\in[0,1].
\]

Thus $G$ is itself a convex minorant of $F$, and it dominates every convex minorant of $F$.
Therefore $G=\operatorname{GCM}(F)$.
\end{proof}

\end{document}
